\section{Algorithm}\label{sec:algorithm}

\subsection{Longest common subsequence by dynamic programming}

For sequences $X = x_1 \dots x_n$ and $Y = y_1 \dots y_m$, the length of the longest common subsequence of the prefixes $x_1 \dots x_i$ and $y_1 \dots y_j$ satisfies
\begin{equation}
    L(i, j) = \begin{cases}
        0 & i = 0 \text{ or } j = 0, \\
        1 + L(i - 1, j - 1) & x_i = y_j, \\
        \max\big(L(i - 1, j),\, L(i, j - 1)\big) & x_i \ne y_j.
    \end{cases}
    \label{eq:lcs}
\end{equation}


The recurrence has overlapping subproblems. Each $L(i, j)$ is needed by up
to three other cells, so the naive recursion recomputes the same values exponentially many times. An
optimal solution for the pair $(i, j)$ contains an optimal solution for a smaller pair, which is what
makes the recurrence valid. A table filled row by row therefore costs $\Theta(nm)$
\citep{clrs}. The subsequence itself is recovered by tracing back through the table.

For the two genomes the table has $29\,903 \times 29\,882 \approx 8.9 \times 10^8$ cells, which is
3.6 GB at 4 bytes per cell, so the textbook algorithm cannot be run on a laptop or in Colab. The
length alone needs only two rows at a time, because a row depends only on the previous one. The
traceback, however, needs the whole table, so discarding the table also discards the alignment.

\subsection{Hirschberg: divide and conquer on top of the recurrence}

Hirschberg's algorithm removes the table without giving up the alignment \citep{hirschberg1975}. It
rests on three facts. First, the last row of the table, $L(\text{mid}, \cdot)$, can be computed with
two rows of memory. Second, for the split of $X$ at its midpoint,
$|LCS(X, Y)| = \max_k \big( |LCS(X[:\text{mid}], Y[:k])| + |LCS(X[\text{mid}:], Y[k:])| \big)$, and
the maximising $k$ is found by adding the forward last row of $X[:\text{mid}]$ against $Y$ to the
last row of the reversed $X[\text{mid}:]$ against the reversed $Y$. Third, once $k$ is fixed the two
halves are independent, so they are solved recursively and their results concatenated. Divide is the
split of $X$, while conquer is the recursion on the two sub-rectangles and combine is the concatenation.
Memory is $O(n + m)$ and time stays $O(nm)$, because the rectangles shrink geometrically
(Section~\ref{sec:complexity}).

Subproblems with at most $40\,000$ cells are solved with a full table and a traceback, which is the
base case of the recursion. On the two genomes the recursion produced 255 split points with a
maximum depth of 7 (Figure~\ref{fig:splits} in the appendix). A single-symbol $X$ is handled as its
own base case, so the split index is always at least 1 and every call makes progress.

\subsection{From subsequence to mutations}

Consecutive matched pairs $(i, j)$ and $(i', j')$ delimit an unmatched gap, $X[i+1:i']$ against
$Y[j+1:j']$. Because the LCS is optimal, the two sides of a gap share no symbol at all: a shared
symbol could be matched without disturbing any other pair, and the LCS would grow. A gap with equal length on both sides is therefore a run of substitutions, a gap with one
empty side is an insertion or a deletion, and anything else is a complex change.



Each substitution inside a coding region is translated in both genomes with the standard genetic
code and classified as synonymous or non-synonymous. The reading frames are not assumed. Gene S and
gene M translate exactly to their protein records, and ORF1ab is translated in two frames because of
the ribosomal frameshift at position 13\,468: the ORF1a frame from 266 and the ORF1b frame from
13\,468. Ten protein records are found whole in the first frame and four in the second, which
validates both the frames and the codon table. Once backtracking (below) locates every protein, the
classification is repeated over all coding regions.

\subsection{Supporting techniques}

Naive search tries every alignment and compares symbol by symbol, which is $O(nm)$ in the worst
case. KMP preprocesses the pattern into its failure table, so that after a mismatch the text pointer
never moves backwards \citep{kmp1977}. The Z-function computes, for \texttt{pattern + \# + text}, the
length of the longest prefix match at every position, and a value equal to $m$ is an occurrence. Both
are $O(n + m)$. The improved methods avoid the repeated work of naive search because after a mismatch
they reuse what is already known about the pattern instead of restarting the comparison.

We consider every open reading frame of at least 100 codons in the three forward frames, each
running from an ATG to the next in-frame stop. Picking the largest set of these that do not
overlap is the activity selection problem. The greedy rule sorts candidates by end position and
keeps each one that does not overlap the last one kept. This is optimal (CLRS ch.~16): the
earliest-ending interval can always replace the first interval of an optimal solution without
creating an overlap. The limitation appears in the results: the objective is the
count of intervals, not the coverage of the genome.

The genetic code has up to six codons per amino acid, so a protein of length $L$
has up to $6^L$ encodings. The search builds an encoding codon by codon. The state of a node is the
set of genome positions still consistent with the codons chosen so far, and a choice that leaves that
set empty is pruned together with its whole subtree. Only the first $k = 6$ amino acids are searched
this way. Every surviving position is then verified by translating the genome forward for the whole
protein.

The candidates are $k$-mers built symbol by symbol, so the search space is a tree of
prefixes. A node carries the genome positions whose window still agrees with the prefix within $d$
mismatches, and the number of those positions is a valid upper bound on the count of any completion
of the prefix, because a window within $d$ mismatches of the full $k$-mer is within $d$ mismatches of
its prefix. Nodes are expanded best first from a max-heap ordered by that bound, the incumbent comes
from a greedy descent, and a node is pruned when its bound cannot beat the incumbent. A hash table is
not used because with $d = 0$ a single pass over the genome solves the problem in $O(nk)$, which is
why the mismatch version is the one that justifies branch and bound.

Centre expansion tries the $2n - 1$ possible centres and is $O(n^2)$ in the worst case. Manacher's
algorithm interleaves a separator so that every palindrome has odd length, then reuses the radius of
the mirror centre inside the rightmost known palindrome \citep{manacher1975}. Each symbol extends the
right boundary at most once, so the algorithm is $O(n)$.
